\chapter{Функция перехода \texorpdfstring{$g$}{g}}
\label{ch:evolution}

\section{Каркас}

\begin{definition}[Локальный CA]
Дискретность $\Lambda$, $(\hl,\htick)$ --- \cref{thm:discreteness,ch:carrier}.
\[
  \Psi^{t+1}=g(\Psi^t),\qquad g=g_{\mathrm{nl}}\circ g_{\mathrm{lin}},
\]
где $g_{\mathrm{lin}}$ --- линейное смешивание по $N(x)$, $g_{\mathrm{nl}}:z\mapsto z\cdot e^{i\Phi(z)}$.
\end{definition}

\begin{definition}[Дискретный лапласиан]
$(\Delta_4 z)(x)=\sum_{y\in N(x)}(z(y)-z(x))$.
\end{definition}

На каркасе не фиксируются $\gamma$, форма $\Phi$, режим $L$ (unitary vs diffusive), $\alpha^*$, $\varepsilon$ ---
они закрываются в \cref{sec:g-law} и \cref{ch:si-sm}.

\section{Ограничения \texorpdfstring{$A_1$--$A_{16}$}{A1-A16}}

\begin{proposition}[Следствия для $g$]
A3 запрещает рост $\sum|z|^2$ --- только bond-unitaries и $\Phi\in\mathbb{R}$.
A4 --- нелинейность только $e^{i\Phi}$.
A5 --- при $|z|\to 0$ gate $\Phi\to\mathrm{large}$.
A6 --- mixing при фиксированной норме.
A7 --- saturating $\Phi$ с $+\varepsilon$.
A8 --- затухание нелинейности при больших $|z|$.
A9 --- CR-метрика holonomy $\zeta$.
\end{proposition}

Линейное $z+\gamma\Delta_4 z$ не унитарно; допустим product локальных unitaries (\cref{sec:g-law}).

\section{Фазовый carrier \texorpdfstring{$\zeta$}{zeta}}

\begin{proposition}
$\Arg(\sum z/z)$ и разность $\atan2$ дают ложный скачок $2\pi$.
Holonomy
\[
  \zeta=(\textstyle\sum_N z)\cdot z^*,\qquad \Delta\phi=\wrap(\Arg\zeta)
\]
--- вход saturating kick (\cref{sec:g-law}).
\end{proposition}

\section{Спин \texorpdfstring{$\mathrm{SU}(2)$}{SU(2)}}

\begin{definition}[Спинор]
$z(x,t)=(z_1,z_2)^\top\in\mathbb{C}^2$, $|z|^2=z^\dagger z$.
\end{definition}

U(1)-фаза $z\to z e^{i\Phi}$ недостаточна для фермионов: $2\pi\to +1$ на $\mathbb{C}$,
но spin-$1/2$ требует $2\pi\to -1$, $4\pi\to +1$.

\begin{definition}[Дискретная rotation gate]
$N_{\mathrm{ring}}=512$, $\omega=e^{i2\pi/N_{\mathrm{ring}}}$, $\Phi\in\mathbb{Z}_{N_{\mathrm{ring}}}$:
\[
  R(\Phi):\ (U,V)\mapsto (U c_\Phi - V s_\Phi,\ U s_\Phi + V c_\Phi)\pmod N,
\]
где $(c_\Phi,s_\Phi)$ --- cos/sin дискретного угла mod $N_{\mathrm{ring}}$.
\end{definition}

\begin{corollary}[Spin-$1/2$]
$R(N_{\mathrm{ring}}/2)\to -z$ на pra-vortex $\hv$ $\Rightarrow$ $U(2\pi)=-\mathbb{1}$, $U(4\pi)=+\mathbb{1}$.
\end{corollary}

Паули: два параллельных спинора в одном $\PlanckCell$ $\Rightarrow$ $z_A+z_B\to 0$;
$K_P$ --- отталкивание при $d\le\hl$.

\section{U(1)$_{\mathrm{vac}}$ и хиральность}

Глобальная фаза $z\to z e^{i\theta}$; holonomy $\zeta$ и Bloch-ось $n=z^\dagger\sigma z/|z|^2$ --- инварианты.
Проекторы $P_L z=(z_1,0)^\top$, $P_R z=(0,z_2)^\top$; в фазе STREAM независимы,
в COLLISION --- смешиваются через $K_P$ и Arg.

\section{Закон \texorpdfstring{$g$}{g} на \texorpdfstring{$\mathbb{Z}_N[i]$}{Z\_N[i]}}
\label{sec:g-law}

\begin{theorem}[Замыкание $g$]
\label{thm:g-law}
Закон $g$ на $Z_{N_{\mathrm{ring}}}[i]$:
\begin{enumerate}
\item leapfrog второго порядка
\[
  Z(x,t+\Delta t)+Z(x,t-\Delta t)=2Z(x,t)+\lfloor\mathcal{N}\rfloor;
\]
\item holonomy $\zeta_{\mathrm{real}},\zeta_{\mathrm{imag}}$ из соседней суммы без $\atan2$;
\item saturating kick
\[
  \Phi=\Bigl\lfloor\frac{K_P\zeta_{\mathrm{imag}}}{\zeta_{\mathrm{real}}+|Z|^2+K_P}\Bigr\rfloor,\quad
  |\Phi|\ge\Delta\phi_{\mathrm{disc}}\ \text{или}\ \Phi=0;
\]
\item $\lfloor\mathcal{N}\rfloor=R(\Phi)-Z$ на обеих компонентах; обратимость A13.
\end{enumerate}
\end{theorem}

\begin{proof}[Идея доказательства]
Каждый пункт --- прямое следствие строк таблицы A1--A16 (\cref{ch:axioms}):
A3+A13 $\Rightarrow$ leapfrog на $\mathbb{Z}$;
A9 $\Rightarrow$ holonomy $\zeta$;
A5+A7 $\Rightarrow$ saturating знаменатель;
A4+A16 $\Rightarrow$ $R(\Phi)$ на $\mathbb{C}^2$.
Congruence ladder CL-1--CL-8 (\cref{sec:congruence}) проверяет согласованность mod $N_{\mathrm{ring}}=512$.
\end{proof}

\subsection{Bit budget \texorpdfstring{$\hv$}{hV}}

\[
  B_{\hv}=\frac{2\pi}{\ln 2}\approx 9.065\ \text{bit},\quad
  N_{\mathrm{ring}}=2^{\lfloor B_{\hv}\rfloor}=512,\quad
  N_\phi=\lceil 4\pi\rceil=13,\quad
  \Delta\phi_{\mathrm{disc}}=41.
\]

\subsection{Congruence ladder}
\label{sec:congruence}

\begin{proposition}[CL-1 -- CL-8]
Mod $N_{\mathrm{ring}}$:
\begin{itemize}
\item CL-1: $Z^++Z^-=2Z+\lfloor\mathcal{N}\rfloor$ $\Leftrightarrow$ $T^{-1}$ (A13);
\item CL-2: $R(N/2)\to -z$ (spin-$1/2$);
\item CL-3: $|\Phi|\ge\Delta\phi_{\mathrm{disc}}$ или $\Phi=0$ (A16);
\item CL-4: $n_E=\lfloor|\Phi|/\Delta\phi_{\mathrm{disc}}\rfloor$, $E=n_E E_0$;
\item CL-5--CL-8: gcd-условия для генератора $\mathbb{Z}_N$, секторов $N_\phi$, ledgers энергии и импульса.
\end{itemize}
\end{proposition}

Заряд $n\in\mathbb{Z}$ (A10) --- не mod $N_{\mathrm{ring}}$.
